\chapter{The meal in the bread}\label{connection:meals}
\regionline{Utensil, wrapper, plate and occasion}
\atlasfiguretitle{The meal in the bread: paired comparison}
\comparisonplate{Rugbrod}{Coco}{A slice for an open sandwich and a fold made to receive a filling.}

Bread is sometimes the thing beside a meal, and sometimes the thing that lets the meal happen. A piece of injera can lift food from the plate. A split arepa can contain a filling. A slice of rugbrød can hold the ingredients of an open sandwich. The bread's shape matters because it changes how food reaches the hand and mouth.

These are useful connections to follow without treating every bread as a substitute for every other. Injera's flexible, porous surface meets sauce differently from the close crumb of a rye slice. A filled pocket and an open sandwich also offer different balances between bread and topping. Put them beside one another and the question becomes what the meal asks the bread to hold.

\section{The hand and the shared dish}
The injera story follows a bread that can form a base and provide pieces for lifting food. Kisra introduces another thin, fermented bread eaten with accompanying dishes. Their batters and methods differ, but both encourage us to look at the meeting between a cooked sheet and the meal around it.\citep{foodishinjera,moggakisra}

Moroccan khobz offers a thicker piece for tearing and dipping. Msemen can move from a layered bread to shredded pieces beneath benrfissa. The dish changes the scale of the bread: a square cooked on the pan becomes smaller pieces within a larger meal. The bread still contributes texture, but it no longer arrives as an individual portion waiting to be eaten whole.\citep{benkhobz,benrfissa}

An illustration of the meal is valuable here because a loaf portrait cannot show every use. Look back at the injera table or the appam plate in the regional chapters. The food around the bread explains why its surface and proportions are worth the trouble. Equipment pictures show how it is made. Meal pictures show what the making is for.

\section{Open, enclosed and folded}
Rugbrød offers a stable base for an open sandwich. The slice remains visible beneath its topping and can be portioned before the meal. Shokupan offers a soft slice for toast and sandwiches, including the meal comparisons in its chapter. Neither needs a dramatic silhouette to have a precise role at the table. Crumb, thickness and the cut surface do the work.\citep{dkrye,namishokupanmeals07}

Bánh mì takes a crusty roll into a filled meal whose ingredients meet along a cut. Andrea Nguyen's account links the bread's Vietnamese history to remembered changes in the filling and manner of eating. A bread introduced through colonial life acquires a meal that needs to be understood in its later setting too. The crust, soft interior, herbs and other fillings make sense together.\citep{nguyenhistorykqed07,nguyenmaster07}

Coco bread gives the filling another kind of opening. Its buttered fold is prepared before baking, then opened to receive a patty. Ftira is cut after baking. Pita's walls separate in the oven, while arepa is opened through its crumb. The meal begins with a space, but the baker creates that space in several different ways.\citep{harriscoco08,ftirafilling}

\section{A familiar bread across an ocean}
Fijian roti brings meals and migration into the same view. The entry places present-day cooking alongside indenture records covering the movement of people from India to Fiji. The National Library of Australia's guide directs descendants to names and places of origin in the emigration passes. That documented movement gives the shared bread name a human route across distance.\citep{nlagirmit09,unescogirmit09}

The home recipe uses a hot-water dough and a pan, with cooked roti kept covered as the stack grows. Compare it with the atta chapati recipe and the difference begins at the bowl. Each can accompany a meal, but a cook needs to follow the particular flour, water and handling rather than assuming that the name guarantees the same dough.\citep{manaliroti06,jangroti09}

Roti also belongs in the langar account of a meal open to all. The Sikh Coalition's explanation places breadmaking within volunteer labour and a shared table that challenges caste barriers. The significance is in the invitation and the work of serving, as well as in the bread itself. A simple disc can be part of a substantial social practice.\citep{sikhlangar06}

\section{Some breads ask us to remember}
Pan de muerto belongs beside these everyday meals without becoming merely another sweet breakfast loaf. Its chapter follows bread into Día de los Muertos observance and the ofrenda. Shape, fragrance and the occasion work together. The eater encounters more than a formula for enriched dough.\citep{conacultamuerto08}

Frybread's story likewise needs its historical setting. The entry connects the bread with forced displacement and ration foods, then follows its place in living communities and debates. A warm, well-made piece can be enjoyed while its history remains present. Delight and difficult memory can belong to the same food.\citep{longwalk08,mihesuahfry08}

The atlas's meals therefore bring several kinds of connection into view. A bread can carry sauce, shelter a filling, make a shared table possible or mark remembrance. Trying an unfamiliar recipe is a way to approach that range with the hands. The most rewarding question after the bake may be what else needs to be on the table for the bread to feel at home.

\section{A meal teaches a shape}
The Sri Lankan pol-roti account offers red curries and lunu miris beside the coconut bread, while the food guide also describes sweeter companions. Another roti preparation, godhamba roti, enters kottu as chopped pieces cooked with vegetables, egg or meat and seasoning. The change is in the meal as well as the bread.\citep{roshanipol06,srilankafoods06,srilankakottu06}

These are satisfying questions to take into a first bake. How large should a torn piece be? Is the sauce meant to meet a crisp edge, a soft interior or several small pieces? The answer can make a familiar shape feel new, and help an unfamiliar bread find a comfortable place at your table.

\begin{breadnote}{Build a meal around the bake}
Choose khobz on page~95, injera on page~85 or bánh mì rolls on page~136. Read the serving account before planning the accompaniments. Notice how the bread is divided and whether it carries, gathers or sits beneath the other food.
\end{breadnote}
